\section{Source audit, proof status, and further research}
\label{audit:section}

\subsection{What has changed relative to the source questions}

The baseline for this continuation is the repository revision
\[
\texttt{\sourcecommit},
\]
dated 11 October 2026. The two paths named
in the research request were inspected at that revision. Six current
incoming archives were extracted for comparison, and relevant canonical
chapters, the manuscript overview, and its proof-status records were
consulted. The source manifest in the package records their paths and
hashes. This fixes what \emph{existing} means in the comparison; it does
not imply that all historical drafts were audited.

The main advancements have different logical forms and should be
integrated accordingly.
\begin{enumerate}
\item The higher-power request for differential-polynomial finite parts
receives an all-polynomial theorem, Theorems~\ref{poly:ML},
\ref{poly:filter}, and \ref{poly:bridge}. The new scalar projection
calculation identifies exactly which polynomial part produces the
nondecaying residue and its Stieltjes correction. Mixed monomials have
a finite, individually convergent word reduction by
Theorem~\ref{poly:words}. The scalar power case is recovered prior work.

\item The canonical report's question about integration by parts at all
weights has the exact direct-sum solution in
Theorem~\ref{poly:IBP}. The endpoint map is retained in
\eqref{poly:IBPintegral}. This solves the formal quotient problem;
classification of the numerical residue-coordinate image still requires
relations among the resulting convergent spectral derivatives.

\item The scalar Gamma--Appell question is sharpened and partly settled
at every degree by Theorem~\ref{ac:main}: the highest height-one Laurent
coefficient is the only \emph{singular} coordinate needed. The explicit
sixth and seventh formulas implement the requested higher-degree
calculations. Elimination of all convergent mixed coordinates is a
different assertion, and remains open.

\item The common-shift problem receives normalized primitives and a
proved endpoint limit in Theorem~\ref{poly:primitivethm}. This establishes
the precise subtraction and its limiting constant. A reduction of every
shifted primitive to a chosen multiple-Hurwitz or Barnes basis is not
proved here.

\item Q8 of the continuous-depth report receives the requested all-order
Gamma--Stieltjes--$Z_\rho$-jet answer in
Theorem~\ref{lj:thm:oddjets}. The reflected jets are supplied with their
own convergent evaluator in Theorem~\ref{lj:thm:Euler}, and the result
extends to ordinary Gamma-weighted polylogarithm derivatives.

\item Conjecture \texttt{conj:firstq} of the canonical fifth-moment
report is proved for all $p=q-2$ by Theorem~\ref{ar:main}.
The general unequal-order conjecture is reduced to an explicit
fixed-gap boundary determinant by Proposition~\ref{ar:gap-det}, but
its nonvanishing is not proved for all larger gaps. The rational
polygamma-product evaluations in Section~\ref{ar:section} include
an independently telescoping example.
\end{enumerate}

\subsection{Corrections and editorial recommendations}

No newly identified false theorem is asserted for the inspected portions
of the six incoming reports. The existing cautions about finite parts,
formal rank, and differentiation of discrete identities are sound.
The following changes are proposed for their integration and for the
research-question lists.

First, update the optimal-prefix conjecture to record the proved
neighboring-order family while retaining the arbitrary-gap statement as a
conjecture. Its determinant and the exact binomial omission convention
should be copied together: retaining the lower-index-zero summand in
\eqref{ar:entries} changes the matrix and invalidates the proof.
Do not replace formal rank by independence of the actual tail
moments or of the even zeta values.

Second, state integration-by-parts reductions as
$\cI(P)=E(Q)+\cI(N)$, with $Q$ at the prescribed weight. For example,
$\cI(u_1^2+u_0u_2)=-\zeta(3)$, not zero. This is an explicit
counterexample to discarding the endpoint constant; the source question
already correctly requires that constant. The new normal-form theorem
provides the requested implementation of that requirement.

Third, descriptions of Gamma--Appell cancellation should distinguish the
singular coordinates from convergent ones. The formulas
\eqref{ac:D6} and \eqref{ac:D7} contain genuine free-order derivatives
of admissible multiple zeta functions. Neither their appearance nor a
failed numerical relation search proves that they are irreducible.
Conversely, calling the singular support a singleton does not state
that the whole answer contains only one nonclassical constant.

Fourth, add the local identity \eqref{lj:eq:local} next to the source's
odd-lattice formula before discussing order derivatives. The known
polynomial values in the discrete index determine the $q=0$ case only.
The reflected derivatives $\partial_s^jZ_\rho(-2m)$ are required for
$q>0$. The same care applies to the normalization factor $2^{1-s}$
in \eqref{lj:eq:Imaster}: its derivatives produce the explicit
$\log 2$ terms in the Stieltjes examples.

Finally, retain the manuscript's existing distinction between its
rejected old Gaussian $S_8$ vector and its revised $S_8$ candidate.
The Gaussian $S_6$ and revised $S_8$ remain conjectural
\cite{Repo}. The symbols $\mathsf S_6,\mathsf S_7$ in this article
are Gamma--Appell finite-part moments, unrelated to those Gaussian
ladder labels. None of the present results changes the Gaussian
proof status.

\subsection{Verification and the role of the accompanying artifacts}
\label{audit:verification}

The proofs in the article establish the all-order statements. The
package also supplies exact finite algebra and independent numerical
diagnostics. Their distinct purposes are recorded in machine-readable
result files.

The adjacent-tail suite checks the source binomial matrix against the
rotated cotangent operator, the positive recurrence against exact
determinants, the full moment reconstruction, the fixed-gap
factorization, and the explicit telescoping certificate.
It records 2604 exact rational assertions and 90 independent symbolic
cotangent-derivative assertions, with adjacent orders through $q=36$.
A deliberately altered coefficient is rejected. Finite examples do not
prove the all-order determinant statement; the positive recurrence does.

The Gamma--Appell suite constructs finite Laurent data, the tail defects,
and the strict word maps. Its weight-six certificate contains 49
shuffle, stuffle, and duality equations on 16 admissible words, together
with an invertible rational subsystem. The separate replay verifies the
certificate and the free-order directional reductions. Integer-index
duality is used for ordinary values only. An additional independent
calculation using ordered set partitions agrees with both printed
directional formulas.

The Lerch suite compares exact Bell coefficients with direct logarithmic
series, checks finite Euler values and trivial zeros, and compares
negative-order Euler sums with positive-order Gamma-moment integrals.
It records 48 exact cumulant checks, 45 finite Euler evaluations
(including 20 odd trivial zeros), two differentiated-product checks,
30 independent mixed-jet numerical comparisons, and direct Stieltjes
derivative checks. The latter comparisons include noninteger and complex
depth. Their largest recorded relative residual is approximately
$1.17\times10^{-71}$.

The differential-polynomial suite verifies all 270 formal monomial
decompositions of weights 2--12, with 22 two-sided exact matrix inverse
identities. It checks 128 mixed and 10 scalar Laurent residues through
weight ten; 48 bilinear pairs, including every applicable zero-order
case; 169 finite harmonic identities; and explicit endpoint constants.
Its numerical companion computes finite parts using a convergent Taylor
integral near zero and ordinary quadrature away from zero. It compares
two splitting points, tests nonlinear endpoint identities, checks
normalized primitives at real and complex arguments, and differentiates
ordinary common-shift integrals.

All numerical records use arbitrary-precision floating-point arithmetic;
they are not interval certificates. The article does not claim
proof-assistant formalization. The symbolic certificates are exact
finite objects, but the analytic continuation, convergence, and
normalization lemmas remain mathematical proofs in the text.

\subsection{Further research questions}
\label{audit:questions}

The following questions seek identities, finite classifications, and
functional relations. Estimates would serve to justify operations or
certificates, rather than being the principal research target.

\subsubsection{The boundary determinant for arbitrary even gaps}

For each fixed $c\ge2$, prove or disprove
$\det F_{b,c}\ne0$ for every $b>c$, with the explicit fundamental
sequences in \eqref{ar:fundamental}. The case $c=2$ is the first
unresolved family. The reduction is useful because the determinant size
depends on the gap, not the upper tail order. A successful proof could
come from a discrete Wronskian identity, a recurrence for the minors,
or a factorization in a suitable polynomial basis. A uniform positivity
argument for all gaps is unavailable: the sign change already recorded
at $c=3$ rules out that simplest claim.

\subsubsection{Systematic rational telescoping for tail products}

The explicit $(4,6)$ identity has a certificate involving only a
polynomial multiple of $A_4(N)A_6(N)$, two single-tail terms, and a
polynomial. Determine whether every adjacent-order reconstruction
identity admits a certificate of this form, and derive its coefficients
directly from the recurrence for $\rho_s$. Such a theorem would replace
spectral continuation by a finite telescoping proof for the whole
ordinary-series family. A useful first output is a bound on the required
polynomial degrees proved from the shift recurrence for the tails.

\subsubsection{The spectral image of the canonical differential basis}

Theorem~\ref{poly:IBP} classifies differential monomials before arithmetic
relations are imposed. Apply the convergent-word map to this basis and
determine its kernel modulo a precisely stated family of free-order
stuffle relations and differential identities. The first interesting
weights are those containing several nonlinear mixed basis elements,
such as $u_0^2u_1^2$ and $u_1^3$ at weight six. The desired result is a
finite presentation with a proved rank in that presentation. Arithmetic
independence of its evaluated numbers should remain a separate question.

\subsubsection{Relations among the sixth and seventh mixed directions}

Determine which combinations in \eqref{ac:D6} and \eqref{ac:D7}
reduce to lower depth, ordinary zeta derivatives, or named
height-one Laurent coefficients. Symmetric free-order stuffle relations
already explain part of the reduction; the surviving antisymmetric
directions require additional identities. Tornheim functional relations,
parameter-dependent iterated integrals, and deformation of known
double-zeta identities are concrete candidates. An identity known only
at integer orders cannot be differentiated for this purpose.

\subsubsection{Higher spectral derivatives of the residue calculus}

The present finite-part map selects the coefficient of $s$ in a
Dirichlet germ. Introduce logarithmically weighted spectral functionals
whose ordinary representatives select $[s^q]$ for $q\ge2$.
Determine whether the same scalar projection controls all their
Stieltjes corrections, and identify a canonical family of normalized
spatial kernels realizing them. For mixed monomials the strict
upper--middle--lower construction already gives convergent higher
partial derivatives; the remaining issue is the exact spatial
normalization and its endpoint constants.

\subsubsection{Rational shifts and multiple-Hurwitz coordinates}

For $I_P(a)$ and $\mathcal P_P(a)$ at rational $a>0$, find finite
formulas in a specified algebra of multiple-Hurwitz derivatives and
cyclotomic polylogarithms. The normalized kernels in
\eqref{poly:primitivekernel} and the endpoint law
\eqref{poly:shiftlimit} fix the constants that such a formula must
respect. Reflection and multiplication combinations may simplify before
individual values do. A complete result should also explain the relation
to the Gamma-quotient normalization with a general base point.

\subsubsection{Additional reductions at rational noninteger depth}

For fixed rational noninteger $\rho>0$, investigate the reflected
coordinates $\partial_\rho^r\partial_s^qZ_\rho(-2m)$.
Theorem~\ref{lj:thm:Euler} already evaluates them by normally convergent
series of finite harmonic-logarithmic expressions. The further
arithmetic question is whether selected depths admit finite expressions
in a predeclared collection of Hurwitz, Gamma, or cyclotomic constants.
Integer-depth termination at $q=r=0$ does not persist automatically
after differentiation. Identifying exactly when an additional finite
reduction exists would be a substantive strengthening of the current
lattice theorem.

\subsubsection{Independent twists and noncentered Gamma weights}

Extend the local lattice construction to the noncentered and tilted
families while keeping order, shift, and depth as independent complex
variables. The centered theorem uses the even Taylor structure and a
single sine factor. Away from the center, the reflection system can
couple two branches or two transforms. Derive the corresponding local
matrix identity and determine its removable values before extracting
mixed jets. The desired output is a finite formula for ordinary
Gamma-weighted root-of-unity polylogarithm derivatives, with every
branch and normalization specified.

\subsubsection{Four or more harmonic tails}

For products of an even number of even-order tails, develop a finite
coordinate system and prove an initial-moment generation theorem.
Centered parity still controls which spectral values can be singular,
but the two-tail rank calculation does not establish the higher-product
claim. The relevant differential operators should contain products of
cotangent kernels and higher pole multiplicities. Exact finite
reconstruction and ordinary telescoping identities would be natural
outputs. Odd numbers of tails require a separate finite-part
normalization because their centered parity changes.

\subsubsection{Functional certificates for the remaining ladders}

Continue the exact functional-equation search for the manuscript's
Gaussian $S_6$ and revised $S_8$ candidates, keeping the evaluated
identity distinct from the formal symbol calculation. The new
convergent spectral coordinates suggest additional Mellin and
polylogarithmic representations, but no bridge to those particular
ladder vectors is proved here. A successful next step should supply an
explicit functional certificate, a proved obstruction in a stated
relation system, or a rigorously nonzero evaluated residual.

\subsubsection{Formal certificates with a small analytic interface}

The finite algorithms admit a useful separation: exact Laurent
extraction and word multiplication, rational matrix reconstruction,
and the analytic lemmas for normal convergence and finite parts.
Formalize the first layer independently, then connect it through proved
interfaces to the scalar residue filter, the normalized
Mittag--Leffler theorem, and the Euler holomorphy theorem.
The purpose is a reusable proof pipeline for further identities,
with each numerical diagnostic remaining a separately identified
experiment.
